\section{Amenazas a la validez}
\paragraph{Validez de constructo.} El uso registrado no es demanda sin restricciones. Las muestras ChronoGraph y las medias Alibaba de intervalos completados tienen semánticas diferentes; la normalización externa no representa bytes o núcleos. Las ausencias sintéticas afectan entradas y dejan objetivos originales para evaluación, sin establecer incidencia o causa de averías reales. La calibración supone prefijos disponibles antes de degradación artificial. Arrastre y edad retienen información anterior a la ventana nominal.

\paragraph{Validez interna.} Estadísticas de prefijo, comprobaciones de partición/huecos, perturbación futura MAR, perturbación de valores ocultos con máscara fija y auditoría de predicciones abordan riesgos identificables. No certifican todo el procesamiento del proveedor. El alcance min--max externo y las fechas de disponibilidad de aristas permanecen inciertos. La validación fuente selecciona hiperparámetros y parada; el destino calibra sin ajustar pesos. Búsqueda estrecha, rotación de condiciones neuronales, ajuste analítico completo de ridge, parametrizaciones diferentes y BatchNorm por microlotes limitan igualdad de oportunidad de optimización. Las enmiendas de memoria y fallos recuperados se conservan.

\paragraph{Validez estadística.} Servicios dependientes, ventanas solapadas y máscaras emparejadas impiden tratar filas como réplicas independientes. Las semillas describen aleatoriedad algorítmica sobre datos fijos. La prueba exacta analiza dirección; la media de diferencias es otro efecto. Reutilizar 101 en selección se estudia solo mediante sensibilidad exploratoria. El bootstrap por bloques requiere aproximación local de dependencia; cambiar 24/48/96 no demuestra estacionariedad universal. Solo cuatro pruebas reciben control principal de multiplicidad. Los demás patrones, horizontes y cohortes son descriptivos.

\paragraph{Validez externa y sesgo histórico.} Fuente y destino ChronoGraph comparten infraestructura y período; la prueba fuente fue inspeccionada históricamente y parte de la prueba actual estuvo expuesta en validaciones anteriores. La selección topológica no es aleatoria. Un protocolo local fechado y exclusión de artefactos localizados no establecen cegamiento externo. Alibaba aporta un día, seis horas de calibración y 36 objetivos para 32 servicios seleccionados por disponibilidad. Sus millones de registros de pods no son millones de objetivos independientes. Normalización y agregación también cambian la tarea. Se requieren otros períodos, organizaciones y regímenes.

\paragraph{Eficiencia, hardware y validez operativa.} El portátil compartido introduce variación; una llamada con tensores residentes no mide desde el colector hasta el controlador. Lote, sincronización, captura e implementación afectan latencia. Los parámetros excluyen estados, preparación y coordinación. La memoria GPU del sistema incluye otras aplicaciones y RSS es muestreado, sin pico continuo. Las épocas excluyen costes de búsqueda y escritura; no identifican un efecto causal exclusivamente arquitectónico. No hay despliegue, controlador, medidor energético fiable ni beneficio SLA medido.

\paragraph{Reproducibilidad.} Los archivos permiten reconstruir cifras y comprobar pesos, pero la imagen derivada depende de una base local preservada. No se acredita una construcción limpia en otra máquina. Las verificaciones del mismo entorno no garantizan igualdad bit a bit entre versiones, dispositivos o kernels. El procesamiento del propietario y las licencias limitan lo que puede redistribuirse. Un repositorio público persistente y sus identificadores siguen pendientes, sin inventar disponibilidad o DOI.
